\let\negmedspace\undefined
\let\negthickspace\undefined
\documentclass[journal]{IEEEtran}
\usepackage[a5paper, margin=10mm, onecolumn]{geometry}
%\usepackage{lmodern} % Ensure lmodern is loaded for pdflatex
\usepackage{tfrupee} % Include tfrupee package

\setlength{\headheight}{1cm} % Set the height of the header box
\setlength{\headsep}{0mm}     % Set the distance between the header box and the top of the text

\usepackage{gvv-book}
\usepackage{gvv}
\usepackage{cite}
\usepackage{amsmath,amssymb,amsfonts,amsthm}
\usepackage{algorithmic}
\usepackage{graphicx}
\usepackage{textcomp}
\usepackage{xcolor}
\usepackage{txfonts}
\usepackage{listings}
\usepackage{enumitem}
\usepackage{mathtools}
\usepackage{gensymb}
\usepackage{comment}
\usepackage[breaklinks=true]{hyperref}
\usepackage{tkz-euclide} 
\usepackage{listings}
% \usepackage{gvv}                                        
\def\inputGnumericTable{}                                 
\usepackage[latin1]{inputenc}                                
\usepackage{color}                                            
\usepackage{array}                                            
\usepackage{longtable}                                       
\usepackage{calc}                                             
\usepackage{multirow}                                         
\usepackage{hhline}                                           
\usepackage{ifthen}                                           
\usepackage{lscape}
\begin{document}

\bibliographystyle{IEEEtran}

\title{8.2.21}
\author{EE25BTECH11054 - S.Harsha Vardhan}
% \maketitle
% \newpage
% \bigskip
\maketitle 
\renewcommand{\thefigure}{\theenumi}
\renewcommand{\thetable}{\theenumi}
\setlength{\intextsep}{10pt} % Space between text and floats

\numberwithin{align}{enumi}
\numberwithin{figure}{enumi}
\renewcommand{\thetable}{\theenumi}
\vspace{-1em}
\textbf{Question 8.2.21}
Find the equation of the conic that satisfies the given conditions:
vertex (0,0) passing through (5,2) symmetric with respect to Y axis.

\solution
The general conic equation is:
\begin{align}
 g(\vec{x}) = \vec{x}^\top \vec{V}\vec{x} + 2\vec{u}^\top \vec{x} + f = 0 \tag{1}
\end{align}
For a parabola with vertex at the origin and axis along the Y-axis, the focus is $\vec{F} = \myvec{0 \\ p}$ and the directrix is $y = -p$. This gives $\vec{n} = \myvec{0 \\ 1}$ and $c = -p$.

\begin{align}
	\vec{V} &= \|\vec{n}\|^2 \vec{I} - \vec{n}\vec{n}^\top 
	= \myvec{1 & 0 \\ 0 & 1} - \myvec{0 \\ 1}\myvec{0 & 1} 
	= \myvec{1 & 0 \\ 0 & 0} \tag{2}
\end{align}
\begin{align}
	\vec{u} &= c\vec{n} - \vec{F} 
	= (-p)\myvec{0 \\ 1} - \myvec{0 \\ p} 
	= \myvec{0 \\ -2p} \tag{3}
\end{align}
\begin{align}
	f &= \|\vec{F}\|^2 - c^2 = p^2 - (-p)^2 = 0 \tag{4}
\end{align}
Substituting these into the general conic equation (1):
\begin{align}
	\vec{x}^\top \myvec{1 & 0 \\ 0 & 0}\vec{x} 
	+ 2\myvec{0 & -2p}^\top\vec{x} + 0 &= 0 \tag{5} \\
	x^2 - 4py &= 0 \tag{6}
\end{align}
The curve passes through the point (5,2).
\begin{align}
	(5)^2 - 4p(2) &= 0 \tag{7} \\
	25 - 8p &= 0 \tag{8} \\
	p &= \frac{25}{8} \tag{9}
\end{align}
The final parameters are:
\begin{align}
	\vec{V} = \myvec{1 & 0 \\ 0 & 0}, \quad
	\vec{u} = \myvec{0 \\ -\frac{25}{4}}, \quad
	f = 0 \tag{10}
\end{align}
The equation of the parabola is:
\begin{align}
	x^2 &= 4\left(\frac{25}{8}\right)y \tag{11} \\
	x^2 &= \frac{25}{2}y \tag{12}
\end{align}
\begin{figure}[h]
    \centering
    \includegraphics[width=1.1\linewidth]{figs/Figure_10.png}
    \caption{}
    \label{fig:placeholder}
\end{figure}

\end{document}